\documentclass{article}

% Language setting
% Replace `english' with e.g. `spanish' to change the document language
\usepackage[english]{babel}

% Set page size and margins
% Replace `letterpaper' with `a4paper' for UK/EU standard size
\usepackage[a4paper,top=2cm,bottom=2cm,left=3cm,right=3cm,marginparwidth=1.75cm]{geometry}

% Useful packages
\usepackage{amsmath}
\usepackage{graphicx}
\usepackage[colorlinks=false]{hyperref}
\usepackage{flafter}



\setcounter{secnumdepth}{4}




\title{Demixing Complex Human Behaviour}

\author{T. Quendera}

\date{\vspace{-5ex}}

\begin{document}
\maketitle

\clearpage





\section*{Author Contributions and Financial Support}
\clearpage

\section*{Acknowledgements}
\clearpage

\begin{abstract}
  an abstract
\end{abstract}
\clearpage

\section*{Thesis Overview}
\clearpage

% table of contents
\tableofcontents
\clearpage

\section*{List of Figures}
\clearpage

\section*{Glossary}
\clearpage


% start content


\section{General Introduction}
  \subsection{Introduction}
  \subsection{Decision Making}
    \subsubsection{How do we make decisions}
    \subsubsection{How do we learn to make them}
      \paragraph{We don’t all learn the same}
  \subsection{Variability}
    \subsubsection{Inter and Intra-Subject variability}
    \subsubsection{Sources of variability}
      \paragraph{Different sources, different signals}
    \subsubsection{The clinic as a fertile ground for studying variability}
    \subsubsection{Discerning a signal amongst noise}
  \subsection{Complexity}
    \subsubsection{What does it mean to be complex}
    \subsubsection{Curse of dimensionality}
    \subsubsection{Tackling complexity}
      \paragraph{Learning from Psychophysics}
      \paragraph{Extending from Psychophysics}
      \paragraph{Modelling complexity}
    \subsubsection{Why studying complexity matters}
      \paragraph{“AI”}
      \paragraph{Do we have the tools and the means?}

\section{Exploring variability in the lab}
  \subsection{Introduction}
    \subsubsection{Decision Making}
        \paragraph{Model Based vs. Model Free learning}
    \subsubsection{Designing experiments across species}
  \subsection{Materials and Methods}
    \subsubsection{Describing the task}
        \paragraph{Foraging Task}
    \subsubsection{Describing the experiments}
        \paragraph{Human vs. Mouse}
        \paragraph{Humans across time}
        \paragraph{Humans across changes in task}
  \subsection{Results}
    \subsubsection{subsubsection name}

  \subsection{Discussion}

Describing the experiments
Humans across time
Humans across changes in task
Describing the tasks
Foraging task
Results
Behaviour is maintained across changes in task
Behaviour is maintained across species
Behaviour is maintained across sessions
Discussion
Acknowledgements
\section{Exploring variability in the clinic}
Introduction
Why translational research on clinical populations
Mental health disorders affect task performance
OCD
Depression
Materials and Methods
Describing the experiments
Depression
OCD
Describing the tasks
Reversal learning task
Stop-signal reaction time task
Foraging task
Results
Reversal learning task
Stop-signal reaction time task
Foraging task
Discussion
Author contributions
Acknowledgements
\section{Exploring variability in the real world}
Introduction
Investigating complex skill learning
Studying behaviour in the wild
Ethologically relevant tasks
Set and setting
Scalability
\section{Materials and Methods}
Hexxed
Task description
Architecture
The game
The database
The analysis pipeline
Optimality in hexxed
ANNs
Other levels
Results
Picky
Sticky
Leapy
Discussion
Author contributions

\section{General Discussion}
Brief summary of the main findings
Opportunities for translation
Concluding remarks








\begin{itemize}
    \item The advent and development of the field of psychophisics drove a lot of behavioural psychologies' findings about human behaviour.
    \item Complex tasks have historically been hard to tackle using the tools of psychophysics.
    \item The capability to clearly highlight significant aspects of simple behaviours is the fields' greatest strength and simultaneously, it's greatest weakness.
    \item The crux of the problem comes from the curse of dimensionality - the harder the problem, the larger the stimulus and action space the sparser the information and our ability to sample from it.
    \item To address and systematically analyze more complex behaviours - playing chess, solving a puzzle, playing an atari videogame - the field has often resorted instead to so-called big data approaches (i.e. using artificial intelligence and machine learning techniques).
    \item When addressing complex problems, some of these agents have proven to be able to solve them at super-human levels (e.g. alpha go)
    \item At the extreme, it is as plausible that they perform better than humans because 1) they are (super)human-like as it is plausible that 2) they are completely non-human.
    \item These ais are a black box - they are not psychophisical in nature
    \item the space of agents grows hyper-exponentially with task complexity.
    \item Performance for performance's sake, than, doesn't mean much. If a deep understanding of how all agents - human or not - is desired, the task must be able to capture fundamental features of performance.
    \item We introduce hexxed - a psychophisics inspired complex behavioral task and large open-access dataset to investigate agents performing complex behaviours.
    \item hexxed is:
    \begin{itemize}
        \item highly parametrizable
        \item deterministic
        \item allows for repeated sampling
        \item highly structured stimulus and action spaces
        which lead to:
        \item comparable across agents
        \item challenging (i.e. hard to solve)
        \begin{itemize}
            \item cognitively... not sensory
        \end{itemize}
    \end{itemize}
\end{itemize}

\section{Methods}

\begin{itemize}
    \item describe hexxed
    \item describe dataset
    \item describe agents
\end{itemize}

\section{Results}
\begin{itemize}
    \item Figures
    \begin{enumerate}
        \item task structure
        \begin{itemize}
            \item printscreen
            \item hypercubes?
            \item mdps?
        \end{itemize}
        \item summary statistics
            \begin{itemize}
            \item how many attempts per level
            \item demographics
            \item how many players per level
            \item players over time
            \end{itemize}
        \item picky
        \item sticky
        \item epiphany
    \end{enumerate}
\end{itemize}

\clearpage

\section*{Figures}

\subsection*{Task introduction}


\subsubsection{lvl 1}

\begin{figure}[h!]
    \centering
    \includegraphics[width=0.8\textwidth]{figures/hex_schematic_1.png}
    \caption{\label{fig:schematic}This frog was uploaded via the file-tree menu.}
\end{figure}

\subsubsection*{rest of levels}



\subsection*{how to solve it and agents}

\subsubsection*{framing the problem}

mdps

\subsubsection*{agents - description}
\subsubsection*{agents - picky}
\subsubsection*{agents - sticky}

\subsection*{Performance/demographics}

\subsection*{humans}

\subsection*{epiphany}


\clearpage

\section{How to collaborate on this file}

\subsection{Adding comments}

Comments can be added by highlighting some text and clicking ``Add comment'' in the top right of the editor pane. To view existing comments, click on the Review menu in the toolbar above. To reply to a comment, click on the Reply button in the lower right corner of the comment. You can close the Review pane by clicking its name on the toolbar when you're done reviewing for the time being.


\subsection{How to include Figures}

First you have to upload the image file from your computer using the upload link in the file-tree menu. Then use the includegraphics command to include it in your document. Use the figure environment and the caption command to add a number and a caption to your figure. See the code for Figure \ref{fig:schematic} in this section for an example.

Note that your figure will automatically be placed in the most appropriate place for it, given the surrounding text and taking into account other figures or tables that may be close by. You can find out more about adding images to your documents in this help article on \href{https://www.overleaf.com/learn/how-to/Including_images_on_Overleaf}{including images on Overleaf}.


\subsection{How to add Tables}

Use the table and tabular environments for basic tables --- see Table~\ref{tab:widgets}, for example. For more information, please see this help article on \href{https://www.overleaf.com/learn/latex/tables}{tables}.

\begin{table}
\centering
\begin{tabular}{l|r}
Item & Quantity \\\hline
Widgets & 42 \\
Gadgets & 13
\end{tabular}
\caption{\label{tab:widgets}An example table.}
\end{table}



\subsection{How to add Lists}

You can make lists with automatic numbering \dots

\begin{enumerate}
\item Like this,
\item and like this.
\end{enumerate}
\dots or bullet points \dots
\begin{itemize}
\item Like this,
\item and like this.
\end{itemize}

\subsection{How to write Mathematics}

\LaTeX{} is great at typesetting mathematics. Let $X_1, X_2, \ldots, X_n$ be a sequence of independent and identically distributed random variables with $\text{E}[X_i] = \mu$ and $\text{Var}[X_i] = \sigma^2 < \infty$, and let
\[S_n = \frac{X_1 + X_2 + \cdots + X_n}{n}
      = \frac{1}{n}\sum_{i}^{n} X_i\]
denote their mean. Then as $n$ approaches infinity, the random variables $\sqrt{n}(S_n - \mu)$ converge in distribution to a normal $\mathcal{N}(0, \sigma^2)$.


\subsection{How to change the margins and paper size}

Usually the template you're using will have the page margins and paper size set correctly for that use-case. For example, if you're using a journal article template provided by the journal publisher, that template will be formatted according to their requirements. In these cases, it's best not to alter the margins directly.

If however you're using a more general template, such as this one, and would like to alter the margins, a common way to do so is via the geometry package. You can find the geometry package loaded in the preamble at the top of this example file, and if you'd like to learn more about how to adjust the settings, please visit this help article on \href{https://www.overleaf.com/learn/latex/page_size_and_margins}{page size and margins}.

\subsection{How to change the document language and spell check settings}

Overleaf supports many different languages, including multiple different languages within one document.

To configure the document language, simply edit the option provided to the babel package in the preamble at the top of this example project. To learn more about the different options, please visit this help article on \href{https://www.overleaf.com/learn/latex/International_language_support}{international language support}.

To change the spell check language, simply open the Overleaf menu at the top left of the editor window, scroll down to the spell check setting, and adjust accordingly.

\subsection{How to add Citations and a References List}

You can simply upload a \verb|.bib| file containing your BibTeX entries, created with a tool such as JabRef. You can then cite entries from it, like this: \cite{greenwade93}. Just remember to specify a bibliography style, as well as the filename of the \verb|.bib|. You can find a \href{https://www.overleaf.com/help/97-how-to-include-a-bibliography-using-bibtex}{video tutorial here} to learn more about BibTeX.

If you have an \href{https://www.overleaf.com/user/subscription/plans}{upgraded account}, you can also import your Mendeley or Zotero library directly as a \verb|.bib| file, via the upload menu in the file-tree.

\subsection{Good luck!}

We hope you find Overleaf useful, and do take a look at our \href{https://www.overleaf.com/learn}{help library} for more tutorials and user guides! Please also let us know if you have any feedback using the Contact Us link at the bottom of the Overleaf menu --- or use the contact form at \url{https://www.overleaf.com/contact}.

\bibliographystyle{alpha}
\bibliography{sample}

\end{document}
